%%% BEGIN fig:dy-abelian-model
\begin{figure}[htbp]
\centering
\begin{tikzpicture}[fd,
  fc lab/.style={font=\small},
  fc note/.style={font=\footnotesize},
  fc amp/.style={hadron blob,ellipse,minimum width=0.95cm,minimum height=1.9cm,inner sep=0pt,font=\footnotesize}]
  % =====================================================================
  % (a) n exchanges between the two eikonal spectator lines
  % =====================================================================
  \begin{scope}
    \node[fc lab] at (-3.15,0) {$\displaystyle\sum_{n\ge0}$};
    \draw[wilson=0.1] (-2.5,0.75) -- (1.9,0.75);
    \draw[wilson=0.1] (-2.5,-0.75) -- (1.9,-0.75);
    \foreach \x in {-1.35,-0.55,1.05} {
      \draw[glauber] (\x,-0.75) -- (\x,0.75);
      \node[vertex] at (\x,0.75) {};  \node[vertex] at (\x,-0.75) {}; }
    \node[fc lab] at (0.25,0) {$\cdots$};
    \node[fc note,anchor=south west] at (-2.5,0.86) {$s_A$：沿$n_A$，荷$g_A$};
    \node[fc note,anchor=north west] at (-2.5,-0.86) {$s_B$：沿$n_B$，荷$g_B$};
    \draw[gray!60!black,line width=0.5pt,{Stealth[length=3.5pt,width=3pt]}-{Stealth[length=3.5pt,width=3pt]}] (2.2,0.7) -- (2.2,-0.7);
    \node[fc lab,anchor=west] at (2.25,0) {$\boldsymbol b$};
    \draw[line width=0.5pt] (-1.45,-1.43) -- (-1.45,-1.55) -- (1.15,-1.55) -- (1.15,-1.43) (-0.15,-1.55) -- (-0.15,-1.65);
    \node[fc note] at (-0.15,-1.93) {$n$次交换（含全部交叉排列）};
    \node[fc lab,anchor=west] at (2.85,0.3) {$\displaystyle=\sum_{n\ge0}\frac{\big(\ii\chi(\boldsymbol b)\big)^n}{n!}=\ee^{\ii\chi(\boldsymbol b)}$};
    \node[fc lab,anchor=west] at (2.85,-0.75) {$\displaystyle\chi(\boldsymbol b)=-\frac{g_Ag_B}{2\pi}\,K_0(m|\boldsymbol b|)$};
    \node[fc lab] at (2.3,-2.6) {(a)\ 程函旁观者之间$n$次有质量阿贝尔玻色子交换之和为相位};
  \end{scope}
  % =====================================================================
  % (b) the two cut placements at O(g^2)
  % =====================================================================
  \begin{scope}[yshift=-4.95cm]
    \begin{scope}[xshift=-1.2cm]
      \draw[wilson=0.76] (-1.6,0.5) -- (1.6,0.5);
      \draw[wilson=0.76] (-1.6,-0.5) -- (1.6,-0.5);
      \node[fc amp] at (-1.6,0) {$\mathcal M^{(0)}$};
      \node[fc amp] at (1.6,0) {$\mathcal M^{(0)*}$};
      \draw[glauber] (-0.62,-0.5) -- (-0.62,0.5);
      \node[vertex] at (-0.62,0.5) {};  \node[vertex] at (-0.62,-0.5) {};
      \draw[cut] (0,-1.15) -- (0,1.15);
      \node[fc note] at (-0.62,0.78) {$s_A$};  \node[fc note] at (-0.62,-0.8) {$s_B$};
      \node[fc lab] at (0,-1.55) {$+\ii\chi(\boldsymbol b)\,|\tilde{\mathcal M}^{(0)}(\boldsymbol b)|^2$};
    \end{scope}
    \node[fc lab] at (1.55,0) {$+$};
    \begin{scope}[xshift=4.3cm]
      \draw[wilson=0.3] (-1.6,0.5) -- (1.6,0.5);
      \draw[wilson=0.3] (-1.6,-0.5) -- (1.6,-0.5);
      \node[fc amp] at (-1.6,0) {$\mathcal M^{(0)}$};
      \node[fc amp] at (1.6,0) {$\mathcal M^{(0)*}$};
      \draw[glauber] (0.62,-0.5) -- (0.62,0.5);
      \node[vertex] at (0.62,0.5) {};  \node[vertex] at (0.62,-0.5) {};
      \draw[cut] (0,-1.15) -- (0,1.15);
      \node[fc note] at (0.62,0.78) {$s_A$};  \node[fc note] at (0.62,-0.8) {$s_B$};
      \node[fc lab] at (0,-1.55) {$-\ii\chi(\boldsymbol b)\,|\tilde{\mathcal M}^{(0)}(\boldsymbol b)|^2$};
    \end{scope}
    \node[fc lab,anchor=west] at (6.55,0) {$=0$};
    \node[fc lab] at (2.3,-2.3) {$\order{g^{2n}}$：$\displaystyle\sum_{j=0}^{n}\frac{(\ii\chi)^j}{j!}\,\frac{(-\ii\chi)^{n-j}}{(n-j)!}=0\ (n\ge1)$；全阶：$|\ee^{\ii\chi(\boldsymbol b)}|^2=1$};
    \node[fc lab] at (2.3,-3.2) {(b)\ 割线位置之和（对旁观者相对横动量$\boldsymbol\Delta$积分，即在$\boldsymbol b$空间中）};
  \end{scope}
\end{tikzpicture}
\caption{阿贝尔模型中旁观者之间的 Glauber 交换（\cref{prop:dy-abelian-model}）。两个旁观者$s_A$、$s_B$近似为沿$n_A$、$n_B$（方向相反）的无限长程函线，荷为$g_A$、$g_B$，横向间距为$\boldsymbol b$；暗红色螺线为质量为$m$的阿贝尔矢量玻色子（Feynman 规范）；时间自左向右，双线上的箭头为荷的流向。
(a)~对$n$次交换（含全部交叉排列）求和给出冲量参数空间的相位：$n$次项为$(\ii\chi)^n/n!$，总和为$\ee^{\ii\chi(\boldsymbol b)}$，$\chi(\boldsymbol b)=-\frac{g_Ag_B}{2\pi}K_0(m|\boldsymbol b|)$为实数（$m\to0$时红外发散，但始终是相位）。
(b)~$\order{g^2}$的两个割图：交换位于割线左侧（振幅）给出$+\ii\chi$，位于右侧（共轭振幅）给出$-\ii\chi$；对旁观者相对横动量$\boldsymbol\Delta$单举的截面中（Parseval 恒等式化为$\boldsymbol b$空间）二者逐阶相消，全阶时$|\ee^{\ii\chi}|^2=1$，但对固定的$\boldsymbol\Delta$逐点并不相消（VF-DY-13，VF-GLA-3）。}
\label{fig:dy-abelian-model}
\end{figure}
%%% END fig:dy-abelian-model

%%% BEGIN fig:dy-oneloop
\begin{figure}[htbp]
\centering
\begin{tikzpicture}[fd,
  t2 lab/.style={font=\footnotesize}]
  % ======================= row 1:  q qbar channel =======================
  % (a) q qbar -> gamma* g, gluon off the quark
  \begin{scope}
    \coordinate (V) at (0,0); \coordinate (E1) at (-2.0,1.15); \coordinate (E2) at (-2.0,-1.15);
    \coordinate (G) at ($(V)!0.45!(E1)$);
    \draw[quark] (E1) -- (G); \draw[quark] (G) -- (V);
    \draw[quark] (V) -- (E2);
    \draw[photon] (V) -- (1.45,0);
    \draw[gluon] (G) -- (0.55,1.2);
    \foreach \p in {V,G} \node[vertex] at (\p) {};
    \path[mom={$p_1$},mom pos=0.45] (E1) -- (G);
    \path[mom'={$p_2$},mom pos=0.6] (E2) -- (V);
    \path[mom'={$q$},mom sep=0.28cm] (V) -- (1.45,0);
    \path[mom={$k$},mom pos=0.55] (G) -- (0.55,1.2);
    \node[t2 lab,anchor=east] at (E1) {$q$};  \node[t2 lab,anchor=east] at (E2) {$\bar q$};
    \node[t2 lab,anchor=west] at (1.45,0) {$\gamma^*$}; \node[t2 lab,anchor=west] at (0.55,1.2) {$g$};
    \node at (-0.3,-1.8) {(a)};
  \end{scope}
  % (b) q qbar -> gamma* g, gluon off the antiquark
  \begin{scope}[xshift=5.0cm]
    \coordinate (V) at (0,0); \coordinate (E1) at (-2.0,1.15); \coordinate (E2) at (-2.0,-1.15);
    \coordinate (G) at ($(V)!0.45!(E2)$);
    \draw[quark] (E1) -- (V);
    \draw[quark] (V) -- (G); \draw[quark] (G) -- (E2);
    \draw[photon] (V) -- (1.45,0);
    \draw[gluon] (G) -- (0.55,-1.2);
    \foreach \p in {V,G} \node[vertex] at (\p) {};
    \path[mom={$p_1$},mom pos=0.6] (E1) -- (V);
    \path[mom'={$p_2$},mom pos=0.45] (E2) -- (G);
    \path[mom'={$q$},mom sep=0.28cm] (V) -- (1.45,0);
    \path[mom'={$k$},mom pos=0.55] (G) -- (0.55,-1.2);
    \node[t2 lab,anchor=east] at (E1) {$q$};  \node[t2 lab,anchor=east] at (E2) {$\bar q$};
    \node[t2 lab,anchor=west] at (1.45,0) {$\gamma^*$}; \node[t2 lab,anchor=west] at (0.55,-1.2) {$g$};
    \node at (-0.3,-1.8) {(b)};
  \end{scope}
  % (c) virtual vertex correction
  \begin{scope}[xshift=10.0cm]
    \coordinate (V) at (0,0); \coordinate (E1) at (-2.0,1.15); \coordinate (E2) at (-2.0,-1.15);
    \coordinate (G1) at ($(V)!0.56!(E1)$); \coordinate (G2) at ($(V)!0.56!(E2)$);
    \draw[quark] (E1) -- (G1); \draw[quark] (G1) -- (V);
    \draw[quark] (V) -- (G2); \draw[quark] (G2) -- (E2);
    \draw[photon] (V) -- (1.45,0);
    \draw[gluon] (G1) -- (G2);
    \foreach \p in {V,G1,G2} \node[vertex] at (\p) {};
    \path[mom={$p_1$},mom pos=0.45] (E1) -- (G1);
    \path[mom'={$p_2$},mom pos=0.45] (E2) -- (G2);
    \path[mom'={$q$},mom sep=0.28cm] (V) -- (1.45,0);
    \node[t2 lab,anchor=east] at (E1) {$q$};  \node[t2 lab,anchor=east] at (E2) {$\bar q$};
    \node[t2 lab,anchor=west] at (1.45,0) {$\gamma^*$};
    \node at (-0.3,-1.8) {(c)};
  \end{scope}
  % ======================= row 2:  q g channel =======================
  % (d) s-channel
  \begin{scope}[xshift=2.3cm,yshift=-4.1cm]
    \coordinate (V1) at (-0.55,0); \coordinate (V2) at (0.65,0);
    \draw[quark] (-2.0,1.05) -- (V1);
    \draw[gluon] (-2.0,-1.05) -- (V1);
    \draw[quark] (V1) -- (V2);
    \draw[photon] (V2) -- (1.95,1.05);
    \draw[quark] (V2) -- (1.95,-1.05);
    \foreach \p in {V1,V2} \node[vertex] at (\p) {};
    \path[mom={$p_1$},mom pos=0.45] (-2.0,1.05) -- (V1);
    \path[mom'={$p_2$},mom pos=0.45] (-2.0,-1.05) -- (V1);
    \path[mom={$q$},mom pos=0.6] (V2) -- (1.95,1.05);
    \path[mom'={$k$},mom pos=0.6] (V2) -- (1.95,-1.05);
    \node[t2 lab,anchor=east] at (-2.0,1.05) {$q$};  \node[t2 lab,anchor=east] at (-2.0,-1.05) {$g$};
    \node[t2 lab,anchor=west] at (1.95,1.05) {$\gamma^*$}; \node[t2 lab,anchor=west] at (1.95,-1.05) {$q$};
    \node at (0,-1.8) {(d)};
  \end{scope}
  % (e) t-channel
  \begin{scope}[xshift=8.0cm,yshift=-4.1cm]
    \coordinate (V1) at (0,0.72); \coordinate (V2) at (0,-0.72);
    \draw[quark] (-1.85,0.72) -- (V1);
    \draw[photon] (V1) -- (1.85,0.72);
    \draw[quark] (V1) -- (V2);
    \draw[gluon] (-1.85,-0.72) -- (V2);
    \draw[quark] (V2) -- (1.85,-0.72);
    \foreach \p in {V1,V2} \node[vertex] at (\p) {};
    \path[mom={$p_1$},mom pos=0.5] (-1.85,0.72) -- (V1);
    \path[mom'={$p_2$},mom pos=0.5] (-1.85,-0.72) -- (V2);
    \path[mom={$q$},mom pos=0.5] (V1) -- (1.85,0.72);
    \path[mom'={$k$},mom pos=0.5] (V2) -- (1.85,-0.72);
    \node[t2 lab,anchor=east] at (-1.85,0.72) {$q$};  \node[t2 lab,anchor=east] at (-1.85,-0.72) {$g$};
    \node[t2 lab,anchor=west] at (1.85,0.72) {$\gamma^*$}; \node[t2 lab,anchor=west] at (1.85,-0.72) {$q$};
    \node at (0,-1.8) {(e)};
  \end{scope}
\end{tikzpicture}
\caption{Drell--Yan 过程 $\order{\as}$ 修正的部分子图（振幅层次，时间自左向右）。上排为 $q\bar q$ 道：(a)、(b)~$q(p_1)\,\bar q(p_2)\to\gamma^*(q)\,g(k)$ 的两个实辐射图，
(c)~单圈顶点修正（与 Born 振幅干涉；无质量外线自能在维数正规化中为零）。下排为 $qg$ 道 $q(p_1)\,g(p_2)\to\gamma^*(q)\,q(k)$ 的两个图：(d)~$s$ 道（内夸克线动量 $p_1+p_2$），(e)~$t$ 道（内夸克线动量 $p_1-q$）
（$\bar q g$ 道的图由电荷共轭得到）。入射反夸克线上的费米子数箭头背离顶点，动量箭头另行标出。}
\label{fig:dy-oneloop}
\end{figure}
%%% END fig:dy-oneloop

%%% BEGIN fig:dy-pole-orientation
\begin{figure}[htbp]
\centering
\begin{tikzpicture}[fd,
  fc hadron/.style={draw=black!70,line width=1.5pt,line cap=butt},
  fc soft/.style={gluon,draw=green!45!black},
  fc state/.style={draw=gray!70!black,line width=0.6pt,dash pattern=on 0.6pt off 1.6pt},
  fc stateB/.style={draw=gray!62,line width=0.6pt,dash pattern=on 0.6pt off 1.6pt},
  fc ax/.style={lcaxis,gray!60!black},
  fc lab/.style={font=\small},
  fc note/.style={font=\footnotesize},
  fc tiny/.style={font=\scriptsize},
  fc pole/.pic={\draw[line width=0.9pt,line cap=round] (-0.075,-0.075)--(0.075,0.075) (-0.075,0.075)--(0.075,-0.075);},
  fc contour/.style={draw=blue!65!black,line width=1.05pt,postaction={decorate,decoration={markings,
      mark=at position #1 with {\arrow{Stealth[length=5pt,width=4.2pt]}}}}},
  fc contour/.default=0.5,
  fc band/.style={fill=red!9,draw=none},
  fc bandA/.style={fill=blue!5,draw=none},
  fc bandB/.style={fill=orange!8,draw=none}]
  % =====================================================================
  % (a) the soft gluon is attached before V_H (tau<0) on both spectators
  % =====================================================================
  \begin{scope}
    % jet subgraphs J_A (top, ordered in x^+) and J_B (bottom, ordered in x^-)
    \fill[fc bandA] (-2.75,0.27) rectangle (2.22,1.95);
    \fill[fc bandB] (-2.75,-1.95) rectangle (2.22,-0.27);
    \node[fc note,text=blue!45!black,anchor=north east] at (2.17,1.5) {$J_A$};
    \node[fc note,text=orange!45!black,anchor=south east] at (2.17,-1.5) {$J_B$};
    % light-cone time axes: x^+ for J_A (top), x^- for J_B (bottom)
    \draw[fc ax] (-3.55,2.75) -- (2.75,2.75) node[right,fc note,text=black] {$x^+$};
    \draw[fc ax] (-3.55,-2.75) -- (2.75,-2.75) node[right,fc note,text=black] {$x^-$};
    \draw[gray!60!black,line width=0.5pt] (0,2.68) -- (0,2.82) (0,-2.68) -- (0,-2.82);
    \node[fc tiny,above] at (0,2.78) {$0$};  \node[fc tiny,below] at (0,-2.78) {$0$};
    \node[fc tiny,above,text=gray!40!black] at (-1.75,2.78) {初态};
    \node[fc tiny,above,text=gray!40!black] at (1.2,2.78) {末态};
    \node[fc tiny,below,text=gray!40!black] at (-1.75,-2.78) {初态};
    \node[fc tiny,below,text=gray!40!black] at (1.2,-2.78) {末态};
    % intermediate states (vertical dotted lines)
    \draw[fc state] (-2.0,1.05) -- (-2.0,1.85);
    \draw[fc state] (-0.7,0.22) -- (-0.7,1.85);
    \draw[fc stateB] (-2.0,-1.05) -- (-2.0,-1.85);
    \draw[fc stateB] (-0.7,-0.22) -- (-0.7,-1.85);
    \node[fc tiny,text=gray!40!black] at (-2.0,2.08) {$D_1$};
    \node[fc tiny] at (-0.15,2.27) {$\dfrac{1}{D_2}=\dfrac{1}{k^-+A+\ii0}$};
    % soft gluon k from s_B to s_A (drawn first: it passes behind a and b)
    \draw[fc soft] (-1.5,-1.6) -- (-1.5,1.6);
    \path[mom'={$k$},mom pos=0.5,mom sep=0.27cm] (-1.5,-1.6) -- (-1.5,1.6);
    % hadrons, splitting vertices, spectators, active partons
    \draw[fc hadron] (-3.6,1.6) -- (-2.5,1.6);
    \draw[fc hadron] (-3.6,-1.6) -- (-2.5,-1.6);
    \draw[plain] (-2.5,1.6) -- (2.75,1.6);
    \draw[plain] (-2.5,-1.6) -- (2.75,-1.6);
    % a and b pass in front of the soft gluon (gap in the background colour of the band)
    \draw[blue!5,line width=3.8pt] (-1.71,1.095) -- (-1.29,0.825);
    \draw[orange!8,line width=3.8pt] (-1.71,-1.095) -- (-1.29,-0.825);
    \draw[quark=0.62] (-2.5,1.6) -- (-0.17,0.11);
    \draw[quark=0.38] (-0.17,-0.11) -- (-2.5,-1.6);
    \draw[photon] (0.2,0) -- (0.9,0);
    \node[hard blob,circle,minimum size=0.4cm,inner sep=0pt] (H) at (0,0) {};
    \foreach \p in {(-2.5,1.6),(-2.5,-1.6),(-1.5,1.6),(-1.5,-1.6)} \node[vertex] at \p {};
    \draw[cut] (2.3,-2.05) -- (2.3,2.05);
    % labels
    \node[fc lab,above] at (-3.25,1.62) {$P_A$};
    \node[fc lab,below] at (-3.25,-1.62) {$P_B$};
    \node[fc lab,right] at (2.75,1.6) {$s_A$};
    \node[fc lab,right] at (2.75,-1.6) {$s_B$};
    \node[fc lab] at (-2.25,0.98) {$a$};
    \node[fc lab] at (-2.25,-0.98) {$b$};
    \node[fc note,anchor=north west] at (0.12,-0.12) {$V_H$};
    \node[fc note,above] at (0.62,0.07) {$\gamma^*$};
    \node[fc lab] at (-0.4,-3.6) {(a)\ $k$在$V_H$之前附着：只进入初态};
  \end{scope}
  % =====================================================================
  % (b) the same Feynman graph, soft gluon attached after V_H (tau>0)
  % =====================================================================
  \begin{scope}[xshift=7.45cm]
    \fill[fc bandA] (-2.75,0.27) rectangle (2.22,1.95);
    \fill[fc bandB] (-2.75,-1.95) rectangle (2.22,-0.27);
    \node[fc note,text=blue!45!black,anchor=north east] at (2.17,1.5) {$J_A$};
    \node[fc note,text=orange!45!black,anchor=south east] at (2.17,-1.5) {$J_B$};
    \draw[fc ax] (-3.55,2.75) -- (2.75,2.75) node[right,fc note,text=black] {$x^+$};
    \draw[fc ax] (-3.55,-2.75) -- (2.75,-2.75) node[right,fc note,text=black] {$x^-$};
    \draw[gray!60!black,line width=0.5pt] (0,2.68) -- (0,2.82) (0,-2.68) -- (0,-2.82);
    \node[fc tiny,above] at (0,2.78) {$0$};  \node[fc tiny,below] at (0,-2.78) {$0$};
    \node[fc tiny,above,text=gray!40!black] at (-1.75,2.78) {初态};
    \node[fc tiny,above,text=gray!40!black] at (1.2,2.78) {末态};
    \node[fc tiny,below,text=gray!40!black] at (-1.75,-2.78) {初态};
    \node[fc tiny,below,text=gray!40!black] at (1.2,-2.78) {末态};
    \draw[fc state] (-1.25,0.58) -- (-1.25,1.85);
    \draw[fc state] (0.62,1.35) -- (0.62,1.85);
    \draw[fc stateB] (-1.25,-0.58) -- (-1.25,-1.85);
    \draw[fc stateB] (0.62,-1.35) -- (0.62,-1.85);
    \node[fc tiny,text=gray!40!black] at (-1.25,2.08) {$D_1$};
    \node[fc tiny] at (0.85,2.27) {$\dfrac{1}{D_2}=\dfrac{-1}{k^--A'-\ii0}$};
    \draw[fc soft] (1.25,-1.6) -- (1.25,1.6);
    \path[mom'={$k$},mom pos=0.5,mom sep=0.27cm] (1.25,-1.6) -- (1.25,1.6);
    \draw[fc hadron] (-3.6,1.6) -- (-2.5,1.6);
    \draw[fc hadron] (-3.6,-1.6) -- (-2.5,-1.6);
    \draw[plain] (-2.5,1.6) -- (2.75,1.6);
    \draw[plain] (-2.5,-1.6) -- (2.75,-1.6);
    \draw[quark=0.7] (-2.5,1.6) -- (-0.17,0.11);
    \draw[quark=0.3] (-0.17,-0.11) -- (-2.5,-1.6);
    \draw[photon] (0.2,0) -- (0.9,0);
    \node[hard blob,circle,minimum size=0.4cm,inner sep=0pt] (H) at (0,0) {};
    \foreach \p in {(-2.5,1.6),(-2.5,-1.6),(1.25,1.6),(1.25,-1.6)} \node[vertex] at \p {};
    \draw[cut] (2.3,-2.05) -- (2.3,2.05);
    \node[fc lab,above] at (-3.25,1.62) {$P_A$};
    \node[fc lab,below] at (-3.25,-1.62) {$P_B$};
    \node[fc lab,right] at (2.75,1.6) {$s_A$};
    \node[fc lab,right] at (2.75,-1.6) {$s_B$};
    \node[fc lab] at (-2.25,0.98) {$a$};
    \node[fc lab] at (-2.25,-0.98) {$b$};
    \node[fc note,anchor=north west] at (0.12,-0.12) {$V_H$};
    \node[fc note,above] at (0.55,0.07) {$\gamma^*$};
    \node[fc lab] at (-0.4,-3.6) {(b)\ $k$在$V_H$之后附着：只进入末态};
  \end{scope}
  % =====================================================================
  % (c), (d): complex k^- plane
  % =====================================================================
  \begin{scope}[shift={(1.2,-5.45)}]
    \fill[fc band] (-0.42,-0.85) rectangle (0.42,0.85);
    \draw[fc ax] (-1.75,0) -- (1.8,0) node[below,fc tiny,text=black,xshift=-3pt] {$\Real k^-$};
    \draw[fc ax] (0,-0.9) -- (0,1.0) node[left,fc tiny,text=black] {$\Imag k^-$};
    \draw[fc contour=0.85] (-1.65,0) -- (1.65,0);
    \pic at (-0.16,-0.26) {fc pole};  \pic at (0.16,0.26) {fc pole};
    \node[fc tiny,anchor=east] at (-0.24,-0.34) {(a)};
    \node[fc tiny,anchor=west] at (0.24,0.34) {(b)};
    \node[fc note,align=center] at (0,-1.52) {(c)\ (a)+(b)：极点分居两侧，\\$k^-$围道被钳住};
  \end{scope}
  \begin{scope}[shift={(5.7,-5.45)}]
    \fill[fc band] (-0.42,-0.85) rectangle (0.42,0.85);
    \draw[fc ax] (-1.75,0) -- (1.8,0) node[below,fc tiny,text=black,xshift=-3pt] {$\Real k^-$};
    \draw[fc ax] (0,-0.9) -- (0,1.0) node[left,fc tiny,text=black] {$\Imag k^-$};
    \draw[fc contour=0.5] (-1.65,0) -- (-1.2,0) .. controls (-0.95,0) and (-0.85,0.62) .. (0,0.62)
        .. controls (0.85,0.62) and (0.95,0) .. (1.2,0) -- (1.65,0);
    \pic at (-0.2,-0.22) {fc pole};  \pic at (0.12,-0.3) {fc pole};
    \node[fc note,align=center] at (0,-1.52) {(d)\ 只经过初态：极点同侧，\\$k^-$围道可移入上半平面};
  \end{scope}
\end{tikzpicture}
\caption{光锥序中间态的取向与$k^-$极点（\cref{lem:dy-poles-orientation}）。Drell--Yan 割图的振幅一侧（割线右侧的共轭振幅为镜像，未画出）按混合光锥序排列：$J_A$（上，浅蓝）按$x^+$、$J_B$（下，浅橙）按$x^-$排序，硬顶点$V_H$在两者中都位于$\tau=0$；
竖直点线标出中间态，$a$、$b$为活动部分子，$s_A$、$s_B$为旁观者，软胶子（绿色）动量$k$由$J_B$流入$J_A$（标准流）。
(a)、(b) 是同一费曼图的两个光锥序项：$k$在$V_H$之前附着于$s_A$时，$k^-$只出现在初态分母$D_2=k^-+A+\ii0$中，零点在下半平面；在$V_H$之后附着时，$k^-$只出现在末态分母$D_2=-(k^--A'-\ii0)$中，零点在上半平面
（$A$、$A'$为与$k^-$无关的实数；$J_B$中对$k^+$有同样的结论，$+\leftrightarrow-$）。取向因而是中间态而不是线的性质：(c)~$k$流经两种取向的中间态时，极点分居实轴两侧，$k^-$围道被钳在 Glauber 区（阴影）内，旁观者之间的交换在$J_A$、$J_B$中同时被钳住，即 Glauber 钳夹；
(d)~$k$只流经初态（例如只附着于$a$）时极点同侧，$k^-$围道可移入上半平面（离开 Glauber 区还须同时形变$k^+$，见\cref{fig:method-glauber-direction}）。}
\label{fig:dy-pole-orientation}
\end{figure}
%%% END fig:dy-pole-orientation

%%% BEGIN fig:dy-two-gluon-exchange
\begin{figure}[htbp]
\centering
\begin{tikzpicture}[fd,
  fc lab/.style={font=\small},
  fc note/.style={font=\footnotesize},
  fc tiny/.style={font=\scriptsize}]
  % =====================================================================
  % (a) box: the two gluons are attached in the same order on both lines
  % =====================================================================
  \begin{scope}
    \draw[gluon] (-0.8,-0.85) -- (-0.8,0.85);
    \draw[gluon] (0.8,-0.85) -- (0.8,0.85);
    \draw[wilson=0.14] (-2.1,0.85) -- (2.1,0.85);
    \draw[wilson=0.14] (-2.1,-0.85) -- (2.1,-0.85);
    \foreach \p in {(-0.8,0.85),(0.8,0.85),(-0.8,-0.85),(0.8,-0.85)} \node[vertex] at \p {};
    \path[mom={$k$},mom sep=0.27cm] (-0.8,-0.85) -- (-0.8,0.85);
    \path[mom'={$K-k$},mom sep=0.27cm] (0.8,-0.85) -- (0.8,0.85);
    \node[fc lab,anchor=east] at (-2.15,0.85) {$v_A,\,\mathbf T_A$};
    \node[fc lab,anchor=east] at (-2.15,-0.85) {$v_B,\,\mathbf T_B$};
    \node[fc tiny] at (0,1.33) {$\dfrac{\ii}{x_A+\ii0}$};
    \node[fc tiny] at (0,-1.35) {$\dfrac{\ii}{-x_B+\ii0}$};
    \node[fc lab] at (0,-2.25) {$J_{\rm box}=\dfrac{2\pi}{ab}\,(\pi+\ii Y)$};
    \node[fc lab] at (0,-3.15) {$c_{\rm box}=X^2$};
    \node[fc lab] at (0,-3.95) {(a)\ 非交叉图（box）};
  \end{scope}
  \node[fc lab] at (3.35,0) {$+$};
  % =====================================================================
  % (b) crossed box
  % =====================================================================
  \begin{scope}[xshift=6.7cm]
    \draw[gluon] (-0.8,-0.85) -- (0.8,0.85);
    % the gluon k passes over the gluon K-k
    \draw[white,line width=7.5pt] (0.2,-0.21) -- (-0.2,0.21);
    \draw[gluon] (0.8,-0.85) -- (-0.8,0.85);
    \draw[wilson=0.14] (-2.1,0.85) -- (2.1,0.85);
    \draw[wilson=0.14] (-2.1,-0.85) -- (2.1,-0.85);
    \foreach \p in {(-0.8,0.85),(0.8,0.85),(-0.8,-0.85),(0.8,-0.85)} \node[vertex] at \p {};
    \path[mom={$k$},mom pos=0.78,mom sep=0.27cm] (0.8,-0.85) -- (-0.8,0.85);
    \path[mom'={$K-k$},mom pos=0.78,mom sep=0.27cm] (-0.8,-0.85) -- (0.8,0.85);
    \node[fc lab,anchor=east] at (-2.15,0.85) {$v_A,\,\mathbf T_A$};
    \node[fc lab,anchor=east] at (-2.15,-0.85) {$v_B,\,\mathbf T_B$};
    \node[fc tiny] at (0,1.33) {$\dfrac{\ii}{x_A+\ii0}$};
    \node[fc tiny] at (0,-1.35) {$\dfrac{\ii}{x_B+\ii0}$};
    \node[fc lab] at (0,-2.25) {$J_{\rm cross}=-\dfrac{2\pi\ii Y}{ab}$};
    \node[fc lab] at (0,-3.15) {$c_{\rm cross}=X^2+\dfrac{C_A}{2}\,X$};
    \node[fc lab] at (0,-3.95) {(b)\ 交叉图（crossed box）};
  \end{scope}
  % =====================================================================
  % sum
  % =====================================================================
  \node[fc lab] at (3.35,-5.15) {$J_{\rm box}\,c_{\rm box}+J_{\rm cross}\,c_{\rm cross}
     =\dfrac{2\pi}{ab}\Big[\underbrace{\pi X^2}_{\mbox{\footnotesize 指数化项}}\;\underbrace{-\,\ii Y\,\tfrac{C_A}{2}\,X}_{\mbox{\footnotesize Regge 相位}}\Big]$};
\end{tikzpicture}
\caption{两条程函线之间的双胶子交换（\cref{prop:dy-eikonal-twogluon}）。两条无限长类时 Wilson 线的速度为$v_A$、$v_B$（快度差$Y>0$，$v_A\cdot v_B=\cosh Y$），色生成元为$\mathbf T_A$、$\mathbf T_B$，时间自左向右；两个胶子取质量$m$（红外正规子）并保留完整传播子，横动量$\mathbf k_T$、$(\mathbf K-\mathbf k)_T$固定，$a=\mathbf k_T^2+m^2$，$b=(\mathbf K-\mathbf k)_T^2+m^2$。
线段旁为两顶点之间的程函传播子（$x_{A,B}=v_{A,B}\cdot k$），其下为纵向积分$J$与色因子$c$（$X=\mathbf T_A\cdot\mathbf T_B$）。
全部$s$道权重在非交叉图中，交叉图只贡献单圈 Regge 相位$\propto C_AY$：横向积分后第一项是指数化项$\frac12(-\ii\chi)^2X^2$，第二项是$\ii\times$实数$\times$厄米色矩阵；因而$p_{\rm opp}=-\ii Y/\pi$，${\rm Re}\,p_{\rm opp}=0$，$\order{g^4}$的幺正性没有残差（VF-REF-9，VF-REF-10）。}
\label{fig:dy-two-gluon-exchange}
\end{figure}
%%% END fig:dy-two-gluon-exchange

%%% BEGIN fig:glauber-exchange
\begin{figure}[htbp]
\centering
\begin{tikzpicture}[fd,
  t2 hadron/.style={draw=black!70,line width=2.2pt,postaction={decorate,decoration={markings,
      mark=at position 0.6 with {\arrow[black!70]{Stealth[length=7pt,width=7.5pt]}}}}},
  t2 lab/.style={font=\footnotesize},
  t2 time/.style={lcaxis,gray!60!black}]
  % ======================= (a) Drell-Yan: initial-state exchange between spectators =======================
  \begin{scope}
    \node[jet blob,circle,minimum size=0.66cm,inner sep=0pt] (A) at (-2.45,1.0) {$A$};
    \node[jetB blob,circle,minimum size=0.66cm,inner sep=0pt] (B) at (-2.45,-1.25) {$B$};
    \node[hard blob,circle,minimum size=0.6cm,inner sep=0pt] (H) at (1.65,1.5) {$H$};
    \draw[t2 hadron] (-3.6,1.0) -- (A.west);
    \draw[t2 hadron] (-3.6,-1.25) -- (B.west);
    % active partons: quark from A, antiquark from B (fermion-number arrow points back into B)
    \draw[quark] (A.12) -- (H.185);
    \draw[quark=0.45] (H.235) -- (B.42);
    \draw[photon] (H.east) -- (2.85,1.5);
    % spectators (remnants) of A and B, continuing into the final state
    \coordinate (SA) at (A.-28); \coordinate (XA) at (2.8,0.25);
    \coordinate (SB) at (B.-8);  \coordinate (XB) at (2.8,-1.2);
    \draw[quark=0.25] (SA) -- (XA);
    \draw[quark=0.25] (SB) -- (XB);
    % Glauber exchange before the hard collision (left of H)
    \coordinate (Gt) at (0.78,2.0); \coordinate (Gb) at (0.78,-2.0);
    \coordinate (GA) at (intersection of SA--XA and Gt--Gb);
    \coordinate (GB) at (intersection of SB--XB and Gt--Gb);
    \draw[glauber] (GA) -- (GB);
    \node[vertex] at (GA) {}; \node[vertex] at (GB) {};
    \node[t2 lab,text=red!55!black,anchor=west] at ($(GA)!0.5!(GB)+(0.16,0)$) {Glauber};
    \node[t2 lab,anchor=west] at (2.85,1.5) {$\gamma^*$};
    \node[t2 lab,anchor=west] at (XA) {$X_A$};
    \node[t2 lab,anchor=west] at (XB) {$X_B$};
    \node[t2 lab,anchor=south] at (-3.25,1.1) {$P_A$};
    \node[t2 lab,anchor=north] at (-3.25,-1.35) {$P_B$};
    \draw[t2 time] (-3.5,-1.95) -- (3.15,-1.95) node[above left,t2 lab,text=gray!50!black] {时间};
    \node at (-0.2,-2.5) {(a)\ DY：旁观者之间的初态相互作用};
  \end{scope}
  % ======================= (b) DIS: final-state exchange between struck quark and remnant =======================
  \begin{scope}[xshift=7.6cm]
    \node[jet blob,circle,minimum size=0.66cm,inner sep=0pt] (A) at (-2.0,-0.95) {$A$};
    \node[hard blob,circle,minimum size=0.6cm,inner sep=0pt] (H) at (0.15,0.45) {$H$};
    \draw[t2 hadron] (-3.5,-0.95) -- (A.west);
    \draw[photon] (-1.35,1.75) -- (H.135);
    \draw[quark] (A.35) -- (H.215);
    \draw[quark] (H.20) -- (2.45,1.45);
    \draw[quark=0.42] (A.-15) -- (2.45,-1.25);
    \coordinate (Gq) at ($(H.20)!0.52!(2.45,1.45)$);
    \coordinate (Gr) at ($(A.-15)!0.735!(2.45,-1.25)$);
    \draw[glauber] (Gq) -- (Gr);
    \node[vertex] at (Gq) {}; \node[vertex] at (Gr) {};
    \path[mom={$q$},mom sep=0.28cm] (-1.35,1.75) -- (H.135);
    \node[t2 lab,anchor=west] at (2.45,1.45) {被击出的夸克};
    \node[t2 lab,anchor=west] at (2.45,-1.25) {靶残余};
    \node[t2 lab,text=red!55!black,anchor=west] at ($(Gq)!0.5!(Gr)+(0.17,0)$) {Glauber};
    \node[t2 lab,anchor=south] at (-3.15,-0.85) {$P$};
    \draw[t2 time] (-3.5,-1.95) -- (2.7,-1.95) node[above left,t2 lab,text=gray!50!black] {时间};
    \node at (-0.45,-2.5) {(b)\ DIS：被击出夸克与靶残余之间的末态相互作用};
  \end{scope}
\end{tikzpicture}
\caption{Glauber 胶子交换（暗红色）。时间自左向右，硬碰撞 $H$ 之前的交换为初态相互作用，之后为末态相互作用。
(a)~Drell--Yan：强子 $A$、$B$ 的旁观者（强子残余，进入末态 $X_A$、$X_B$）在硬碰撞之前交换 Glauber 胶子（$|k^+k^-|\ll k_T^2$）；来自 $B$ 的入射反夸克线与 $A$ 的旁观者线的交叉不是顶点。
在单个割图中，$A$ 一侧（旁观者线与活动部分子线）的 $k^-$ 极点与 $B$ 一侧的 $k^+$ 极点都分居实轴两侧，围道被钳在 Glauber 区（\cref{fig:glauber-contour}），只有对旁观者的末态切割求和（幺正性）之后才能把围道变形出 Glauber 区；
(b)~DIS：被击出的夸克与靶残余在硬碰撞之后交换 Glauber 胶子；对单举 DIS 它们由幺正性相消，而在 SIDIS 的 TMD 中它们组合成未来指向的 Wilson 线，
与 Drell--Yan 中由初态相互作用产生的过去指向 Wilson 线相对照（Sivers 函数因而在两过程间变号）。}
\label{fig:glauber-exchange}
\end{figure}
%%% END fig:glauber-exchange
